\section{Additional results}
\label{app:more}

\subsection{Structural identifiability runs}
\label{app:structural}

The symbolic analyses used Julia and the package StructuralIdentifiability.jl in the versions
\texttt{\juliaversion} and \texttt{\siversion}, on the Bergman-type model with the two-compartment gut cascade in its
active regime and glucose as the output. Seven formulations were run; the four of Table~\ref{tab:structural} are the ones
that differ in outcome, and the raw output of every run is archived with the code. The tool's local assessment
reports all three parameters locally identifiable. Where the gut initial conditions were left free or
unknown the rates were reported locally but not globally identifiable; in canonical coordinates
$\sigma_1=\ke+\ka$ and $c=\ke\ka$ the two combinations were globally identifiable.

\subsection{Tables of the main analyses}

\begin{table}[t]
\centering\small
\begin{tabular}{lccc}
\toprule
Parameter & $0.5\times$ & $1\times$ & $2\times$ \\
\midrule
$S_I$ & 30 (67\%) & 23 (51\%) & 5 (11\%) \\
$k_e$ & 45 (100\%) & 40 (89\%) & 26 (58\%) \\
$k_a$ & 44 (98\%) & 39 (87\%) & 34 (76\%) \\
interior for $S_I$ & 15 & 22 & 40 \\
interior, original definition & 0 & 0 & 0 \\
\bottomrule
\end{tabular}
\caption{Subjects whose maximum-likelihood estimate lies on a bound (iAUC objective, 45 subjects), for three box widths. The $1\times$ box is the primary one.}
\label{tab:bounds}
\end{table}

\begin{table}[t]
\centering\small
\begin{tabular}{lccc}
\toprule
Objective, parameter & $0.5\times$ & $1\times$ & $2\times$ \\
\midrule
iAUC, $S_I$ & 0 [0, 8] & 7 [2, 18] & 53 [39, 67] \\
iAUC, $k_e$ & 0 [0, 8] & 0 [0, 8] & 0 [0, 8] \\
iAUC, $k_a$ & 0 [0, 8] & 0 [0, 8] & 0 [0, 8] \\
\bottomrule
\end{tabular}
\caption{Percentage of subjects with a bounded 95\% profile-likelihood interval (Wilson 95\% interval in brackets). A dash means the analysis was not run.}
\label{tab:profiles}
\end{table}

\begin{table}[t]
\centering\small
\begin{tabular}{lccc}
\toprule
Box, evaluated at & iAUC & iAUC+centroid & trace \\
\midrule
$0.5\times$, own estimate & 2.3e+05 & 4e+05 & 6.4e+03 \\
$0.5\times$, common reference & 2.3e+05 & 2.1e+05 & 1.2e+03 \\
$1\times$, own estimate & 4.3e+05 & 2.8e+07 & 8.8e+07 \\
$1\times$, common reference & 3.3e+05 & 8.7e+04 & 8.1e+02 \\
$2\times$, own estimate & 7.8e+03 & 5e+14 & 1.1e+14 \\
$2\times$, common reference & 7.8e+03 & 1e+05 & 4.1e+02 \\
\bottomrule
\end{tabular}
\caption{Median condition number of the Schur-complement timing block, at each objective's own maximum-likelihood estimate and at one common reference point per subject.}
\label{tab:ladder}
\end{table}

\begin{table}[t]
\centering\small
\begin{tabular}{lccc}
\toprule
Pair & Difference [95\% CI] & Wins & Holm $p$ \\
\midrule
grad3 $-$ grid3 & -1 [-3, 1] & 25/45 & 1.000 \\
grad1 $-$ grid1 & 0 [-1, 1] & 24/45 & 1.000 \\
grad3 $-$ grad1 & -34 [-68, -12] & 25/45 & 0.435 \\
grad3\_trace $-$ grad3 & 32 [2, 67] & 21/45 & 1.000 \\
grad1\_trace $-$ grad1 & 39 [15, 91] & 13/45 & 0.178 \\
grad3 $-$ grid1 & -34 [-68, -11] & 27/45 & 0.435 \\
grad3 $-$ rf & -129 [-220, -54] & 28/45 & 0.093 \\
grad3 $-$ snpe & -395 [-612, -254] & 34/45 & 0.000 \\
\bottomrule
\end{tabular}
\caption{Pre-specified paired comparisons (negative: the first cell predicts better).}
\label{tab:comparisons}
\end{table}

\begin{table}[h]
\centering\small
\begin{tabular}{lcccrr}
\toprule
box & level (\%) & rise $\delta$ & $\SI$ bounded, interior (\%) & $\ke$ bounded (\%) & $\ka$ bounded (\%) \\
\midrule
$0.5\times$ & 90 & 1.35 & 0 / 15 (0) & 0 & 0 \\
$0.5\times$ & 95 & 1.92 & 0 / 15 (0) & 0 & 0 \\
$0.5\times$ & 99 & 3.32 & 0 / 15 (0) & 0 & 0 \\
$1\times$ & 90 & 1.35 & 5 / 22 (23) & 0 & 0 \\
$1\times$ & 95 & 1.92 & 3 / 22 (14) & 0 & 0 \\
$1\times$ & 99 & 3.32 & 0 / 22 (0) & 0 & 0 \\
$2\times$ & 90 & 1.35 & 26 / 40 (65) & 0 & 0 \\
$2\times$ & 95 & 1.92 & 24 / 40 (60) & 0 & 0 \\
$2\times$ & 99 & 3.32 & 19 / 40 (48) & 0 & 0 \\
\bottomrule
\end{tabular}
\caption{Bounded profile intervals under iAUC at three interval levels, recomputed from the stored profiles. The rise $\delta$ is half the corresponding quantile of $\chi^2_1$.}
\label{tab:levels}
\end{table}


\subsection{Noise budgets on the replica}

\begin{table}[h]
\centering\small
\begin{tabular}{llcccc}
\toprule
CGM noise & carbohydrate error & gradient & grid & personal mean & ratio to real \\
\midrule
none & none & 2 & 15 & 1800 & 0.00 (0.00--0.00) \\
none & 25\% & 710 & 707 & 1800 & 0.23 (0.21--0.26) \\
on & none & 2394 & 2391 & 2874 & 0.79 (0.72--0.89) \\
on & 25\% & 2528 & 2531 & 2874 & 0.84 (0.76--0.94) \\
on & 50\% & 2829 & 2829 & 2874 & 0.94 (0.85--1.04) \\
\bottomrule
\end{tabular}
\caption{What-if noise budgets: held-out iAUC error of the replica (mg\,dL$^{-1}$\,min) under each noise setting. The real-data error of the gradient fit is 3025. The ratio carries a bootstrap interval over subjects. The settings are not additive shares of the real error.}
\label{tab:noise}
\end{table}


\subsection{Subject-level correlates}

\begin{table}[h]
\centering\small
\begin{tabular}{lccc}
\toprule
 covariate & $\SI$ interval bounded & timing rate on the upper bound & rise of the $\SI$ profile \\
\midrule
number of meals & -0.04 (0.771; 1.000) & -0.36 (0.015; 0.151) & 0.34 (0.023; 0.207) \\
carbohydrate range & 0.08 (0.622; 1.000) & -0.25 (0.104; 0.753) & 0.25 (0.099; 0.753) \\
mean observed iAUC & 0.06 (0.687; 1.000) & 0.65 (0.000; 0.000) & -0.21 (0.165; 0.823) \\
$\SI$ estimate & 0.10 (0.525; 1.000) & -0.61 (0.000; 0.000) & 0.25 (0.094; 0.753) \\
\bottomrule
\end{tabular}
\caption{Spearman correlation across subjects at the primary box, with the raw and the Holm-adjusted $p$-values in parentheses. Exploratory.}
\label{tab:correlates}
\end{table}


\subsection{Further figures}

\begin{figure}[h]
\centering
\safefigure[width=\linewidth]{a1_gut_and_boxes}
\caption{\textbf{(a)} Largest and median change in simulated glucose under the exchange $\ke\leftrightarrow\ka$
as a function of residual intestinal content, as a percentage of the meal (shaded: $5$ to $10$\,mg\,dL$^{-1}$, for scale).
\textbf{(b)} Subjects whose estimate lies on a bound, by box width.}
\label{fig:gutsweep}
\end{figure}

\begin{figure}[h]
\centering
\safefigure[width=0.55\linewidth]{a2_generic_rank}
\caption{Median scaled singular values of the stacked sensitivity matrix, smallest first, by observable.}
\label{fig:rank}
\end{figure}

\begin{figure}[h]
\centering
\safefigure[width=\linewidth]{a3_dalla_man}
\caption{Second model class (Dalla Man), iAUC objective. \textbf{(a)} Eigenvalues of the Fisher information in
log coordinates, over subjects. \textbf{(b)} Subjects with a bounded $95\%$ profile interval per parameter.}
\label{fig:dm}
\end{figure}

\begin{figure}[h]
\centering
\safefigure[width=0.6\linewidth]{a4_robustness}
\caption{Median Kendall's $\tau$ between the reference ordering of the three parameters by the projected
gradient diagnostic and the ordering under each factor of H10; dotted: the threshold $0.8$.}
\label{fig:robustness}
\end{figure}
